\section{Pruebas}

El proyecto tiene \totalpruebas{} pruebas automáticas, que cubren el
\coberturapruebas{}\,\% de las líneas y ramas del motor. El código pasa además
\texttt{ruff} y \texttt{mypy} en modo estricto. Aparte de las pruebas por caso (estructura
vacía, un solo elemento, división y fusión de nodos, desbordamiento de página), se usaron
tres formas de verificación:

\begin{itemize}[nosep, leftmargin=*]
  \item \emph{Validadores de estructura.} Cada índice tiene una función que lo recorre
        completo y comprueba sus invariantes. En el B+ revisa el orden de las claves, la
        ocupación mínima y que todas las hojas estén al mismo nivel; en el hash, que cada
        bucket tenga $2^{gd-ld}$ punteros; en el R-Tree, que cada MBR sea el mínimo.
  \item \emph{Comparación contra un modelo.} A cada estructura se le aplican miles de
        operaciones aleatorias y el resultado se compara con una lista ordenada o con una
        búsqueda por fuerza bruta. Las consultas espaciales se ejecutan con y sin R-Tree y
        deben devolver lo mismo.
  \item \emph{Comparación con PostGIS.} Las 1\,800 consultas del experimento espacial
        devuelven las mismas filas que PostGIS.
\end{itemize}

Las pruebas de almacenamiento usan páginas de 256 bytes para que con pocos registros ya
haya varias páginas y se ejerciten las divisiones.
